% SPDX-License-Identifier: Apache-2.0
%
% AXI in TxHDL: the tutorial on the bus. Built by //docs:axi. Every
% listing is included from the tree and every printout and figure is
% what the build made by running the examples.
\documentclass[journal,9pt]{IEEEtran}

% Listings of Rust set by rustfmt at 80 columns: 5.6pt is what fits
% the frame of a column (issue 1061).
\newcommand{\listingsize}{\fontsize{5.6}{6.6}\selectfont}
% SPDX-License-Identifier: Apache-2.0
%
% The house preamble, shared by every document in this directory. Loaded
% after \documentclass, before \title. Keeping it in one file is what
% stops two articles drifting apart in font, listing style or figure
% style.
% Computer Modern. IEEEtran selects Times (ptm) by default, so the face has
% to be chosen explicitly.
%
% Latin Modern is the Computer Modern variety used here, and the choice is
% forced rather than aesthetic. Under T1 encoding, plain `cmr` has no Type 1
% outlines unless cm-super is installed, and it is not in the pinned
% distribution, so pdflatex falls back to EC bitmaps and every glyph in the
% PDF becomes a Type 3 font. That was measured: the first build produced a
% document with 10 Type 3 fonts and no Type 1. Latin Modern is Computer
% Modern redrawn with a full T1 repertoire and real .pfb outlines, so the
% same design comes out scalable.
\usepackage[T1]{fontenc}
\usepackage[utf8]{inputenc}
\usepackage{lmodern}
% microtype is not in the pinned TeX distribution. emergencystretch alone
% keeps the two-column measure from producing overfull lines.
\setlength{\emergencystretch}{3em}

\usepackage{amsmath}
\usepackage{amssymb}
\usepackage{amsthm}
\usepackage{booktabs}
\usepackage{array}
% enumitem is not in the pinned TeX distribution; the list spacing below
% does what \setlist would have done.
\usepackage{float}
\usepackage{xcolor}
\usepackage{listings}
\usepackage{tikz}
\usetikzlibrary{arrows.meta,positioning,fit,backgrounds,calc}
\usepackage[hidelinks,bookmarks=true]{hyperref}

% Numbered, definition-style box for a rule the reader is meant to apply.
\theoremstyle{definition}
\newtheorem{rulebox}{Rule}
\newtheorem{example}{Example}

% Unnumbered asides.
\theoremstyle{remark}
\newtheorem*{tip}{Tip}
\newtheorem*{pitfall}{Pitfall}

% Tighter lists, in place of enumitem's \setlist.
\setlength{\topsep}{2pt}
\setlength{\itemsep}{1pt}
\setlength{\parsep}{0pt}

% Code in the running text. The typewriter face under T1 ligates `<<`
% and `>>` into guillemets, so a shift printed as \code{<<} came out as
% a French quotation mark (issue 571). microtype's \DisableLigatures is
% not in the pinned distribution, but the primitive it rests on is
% pdfTeX's own: \pdfnoligatures strips every ligature from the font it
% is given, and the current font inside \texttt is the typewriter face
% at the size in use, so each size is stripped the first time \code
% sets it. Code wants no ligature at all: two quotes are two quotes.
\newcommand{\code}[1]{\texttt{\ifdefined\pdfnoligatures\pdfnoligatures\font\fi #1}}
\newcommand{\term}[1]{\emph{#1}}

% Syntax highlighting. Four roles, distinguishable in colour and still
% distinguishable in greyscale, because an IEEE article gets printed: the
% keyword is the only bold role, the comment is the only italic one, and the
% two remaining roles differ in lightness as well as in hue.
\definecolor{kwblue}{RGB}{20,60,150}
\definecolor{cmtgreen}{RGB}{90,120,90}
\definecolor{strred}{RGB}{150,50,40}
\definecolor{numgray}{gray}{0.45}
\definecolor{framegray}{gray}{0.70}
\definecolor{bgshade}{gray}{0.97}
% A darker fill for figure cells. The listing background has to stay very
% light to keep code readable; a figure cell has to be clearly shaded or the
% distinction it draws is invisible in print.
\definecolor{cellshade}{gray}{0.90}

% The listing size is the document's choice, because it depends on the
% column: a two-column article sets `\listingsize` to 6pt and keeps its
% listings within 70 columns, or to 5.6pt when it includes source that
% rustfmt sets at 80, which is what fits a column's frame (issue 1061);
% a one-column document keeps the body size and includes source files
% at 80. `//docs:frame_test` checks every listing line against its frame. Lines are never broken; a line that does not fit
% overflows the frame, which is visible, rather than wrapping, which is
% not.
\providecommand{\listingsize}{\scriptsize}

% `'` is deliberately not a string delimiter: the language uses it for loop
% labels, as Rust does, so treating it as a quote would swallow the rest of
% the line.
\lstdefinestyle{txhdl}{
  basicstyle=\ttfamily\listingsize,
  keywordstyle=\bfseries\color{kwblue},
  commentstyle=\itshape\color{cmtgreen},
  stringstyle=\color{strred},
  numberstyle=\tiny\color{numgray},
  morekeywords={mod,use,pub,const,type,struct,enum,trait,impl,fn,async,let,mut,
    match,if,else,loop,while,for,in,break,continue,return,as,await,
    module,bus,channel,transaction,protocol,pipeline,flow,seq,config,
    port,stage,pipe,instance,connect,bind,tag,domain,blackbox,foreign,
    property,role,signal,handler,true,false,
    sequence,interface,modport,implements,package,import,var,wire,func,proc,
    case,default,next,fork,branch,join,state,fsm},
  morecomment=[l]{//},
  morestring=[b]",
  showstringspaces=false,
  columns=fullflexible,
  keepspaces=true,
  breaklines=false,
  % Framed, so a listing is bounded rather than floating in the column.
  frame=single,
  rulecolor=\color{framegray},
  framesep=3pt,
  backgroundcolor=\color{bgshade},
  xleftmargin=3pt,
  xrightmargin=3pt,
  aboveskip=6pt,
  belowskip=6pt,
  % Every listing is numbered and captioned, so it can be referred to by
  % number. `caption` without `float` numbers the listing and leaves it
  % where it was written, which is what a listing in running prose needs.
  captionpos=b,
  abovecaptionskip=2pt,
  belowcaptionskip=6pt,
}

% Figures drawn as text. Framed the same way, and with the highlighting
% switched off, because a box-and-arrow diagram is not code and half its
% words turning blue reads as an accident. Setting a style adds keys rather
% than clearing the ones already set, so the three roles are neutralised by
% hand rather than by leaving them out. Program output is printed in this
% style too, and a netlist's assignment can run past the frame, so a line
% that does not fit is broken and the rest indented; source never is.
\lstdefinestyle{ascii}{
  basicstyle=\ttfamily\listingsize,
  keywordstyle=\ttfamily,
  commentstyle=\ttfamily,
  stringstyle=\ttfamily,
  columns=fullflexible,
  keepspaces=true,
  breaklines=true,
  breakindent=2em,
  frame=single,
  rulecolor=\color{framegray},
  framesep=3pt,
  backgroundcolor=\color{bgshade},
  xleftmargin=3pt,
  xrightmargin=3pt,
  aboveskip=6pt,
  belowskip=6pt,
}
\lstset{style=txhdl}

% House TikZ styles. `shorten` lives on the arrow styles rather than on each
% arrow, so every arrow in the document gets the gap, including ones added
% later. TikZ clips a path at a node's border, so without it an arrowhead
% butts against the box with no gap at all. Keep `node distance` well above
% twice the shorten, or the arrow shrinks to nothing.
\tikzset{
  box/.style={draw,rounded corners=2pt,align=center,inner sep=4pt,
              font=\footnotesize},
  boxf/.style={box,fill=bgshade},
  lbl/.style={font=\scriptsize\itshape,align=center},
  fl/.style={-{Stealth[length=4pt]},shorten >=2.5pt,shorten <=2.5pt},
  fld/.style={fl,dashed},
}


% The contents list: the class gives a section's number 2.75em, which
% a Roman numeral past thirty overruns onto the title, so the box is
% wider here, and a subsection's entry starts where a title does.
\makeatletter
\def\l@section#1#2{\addpenalty{\@secpenalty}\addvspace{1.0em plus 1pt}%
    \@tempdima 5.75em \begingroup \parindent \z@ \rightskip \@pnumwidth%
    \parfillskip-\@pnumwidth {\bfseries\leavevmode #1}\hfil\hbox to\@pnumwidth{\hss #2}\par%
    \endgroup}
\def\l@subsection{\@dottedtocline{2}{5.75em}{5em}}
\def\l@subsubsection{\@dottedtocline{3}{10.75em}{5.5em}}
\makeatother


\InputIfFileExists{buildstamp.tex}{}{}
\providecommand{\buildstamp}{Draft build (outside the Bazel flow).}

% A range of a simulation-only source, by its begin/end markers.
\newcommand{\simpart}[3]{%
\lstinputlisting[rangebeginprefix=//\ begin\{,rangebeginsuffix=\},%
rangeendprefix=//\ end\{,rangeendsuffix=\},includerangemarker=false,%
linerange=#1-#1,caption={#2},label={#3}]{lib/parts/src/bus/axi/sim.rs}}

% A range of the AXI-Lite bridge, by its begin/end markers.
\newcommand{\litepart}[3]{%
\lstinputlisting[rangebeginprefix=//\ begin\{,rangebeginsuffix=\},%
rangeendprefix=//\ end\{,rangeendsuffix=\},includerangemarker=false,%
linerange=#1-#1,caption={#2},label={#3}]{lib/parts/src/bus/axi_lite.rs}}


% A range of the simulation-only Wishbone parts.
\newcommand{\simwbpart}[3]{%
\lstinputlisting[rangebeginprefix=//\ begin\{,rangebeginsuffix=\},%
rangeendprefix=//\ end\{,rangeendsuffix=\},includerangemarker=false,%
linerange=#1-#1,caption={#2},label={#3}]{lib/parts/src/bus/wb/sim.rs}}

% A range of the fuzzing of the link, by its begin/end markers.
\newcommand{\fuzzpart}[3]{%
\lstinputlisting[rangebeginprefix=//\ begin\{,rangebeginsuffix=\},%
rangeendprefix=//\ end\{,rangeendsuffix=\},includerangemarker=false,%
linerange=#1-#1,caption={#2},label={#3}]{lib/parts/src/bus/axi/fuzz.rs}}

% A range of the library's source, by its begin/end markers.
\newcommand{\axipart}[3]{%
\lstinputlisting[rangebeginprefix=//\ begin\{,rangebeginsuffix=\},%
rangeendprefix=//\ end\{,rangeendsuffix=\},includerangemarker=false,%
linerange=#1-#1,caption={#2},label={#3}]{lib/parts/src/bus/axi.rs}}

\title{AXI in TxHDL}

\author{Filip Filmar and Dragiša Janković%
\thanks{Both authors are independent. Filip Filmar is the
corresponding author, at f@hdlfactory.com; Dragiša Janković is
reached at linkedin.com/in/dragisa-jankovic.}%
\thanks{This is the tutorial on one bus of the language's library,
for a reader who has met TxHDL through the step-by-step
tutorial~\cite{tutorial} and wants to put an AXI link in a design,
or to read how the link is made. Every listing is included from the
project repository \code{HDL/txhdl} at build time and every printout
and figure is what the build produced by running the examples. The
reservation station has a tutorial of its own~\cite{station}, the
runtime is in the runtime document~\cite{runtime}, and the
cover~\cite{cover} lists every document. The AXI examples are here
and not in the examples document~\cite{examples}, because a bus is a
subsystem and not one more construct.}%
\thanks{The authors used a large language model, Claude, as an
assistant in exploring the concepts and in writing and constructing
the documents and the programs; every commit of the repository says
so and records the prompts that led to it, verbatim.}%
\thanks{\buildstamp}}

\markboth{AXI in TxHDL}{Filmar and Janković: AXI in TxHDL}

\begin{document}
\maketitle

\begin{abstract}
A hardware channel transmits one transaction in one direction. AXI4
coordinates five concurrent channels. A client interacting with those
channels directly must assemble write addresses with data beats, track
lengths, inspect \code{last}, allocate transaction identifiers, and reorder
completions. TxHDL abstracts these mechanics through the \code{axi()}
constructor. The function constructs the complete AXI4 channel bundle and
yields two typed endpoints, mirroring the pair returned by \code{chan()}. A
host client submits a burst and awaits completion. A peripheral client
accepts fully decoded transactions and returns responses from concurrent
processes in arbitrary order. Neither endpoint exposes individual channel
signals directly. Two protocol engines connect these endpoints. The engines
lower to synthesizable Verilog and VHDL, verified against cycle-accurate
simulation traces. An address-mapped router multiplexes one host link
across multiple peripherals, returning decode errors locally when addresses
miss. Inserting the router requires only channel reconnection without
changing client logic. For simpler peripherals, TxHDL provides AXI-Lite
endpoints without identifiers or bursts. A synthesized bridge maps AXI4
bursts into single-beat AXI-Lite transactions across peripheral banks.
\end{abstract}

\section{What a Link Is}
\label{sec:a-what}

The AXI4 protocol specifies five independent channels. Write address,
write data, and write response operate over \code{aw}, \code{w}, and
\code{b}. Read address and read data operate over \code{ar} and \code{r}. A
write transaction consists of an address beat, one or more data beats
terminated by an asserted \code{last} flag, and a completion response. A
read transaction pairs an address beat with a sequence of data beats.
Each beat holds a transaction identifier. Interleaving beats with distinct
identifiers allows multiple concurrent transactions in flight.

TxHDL shields client code from channel-level bookkeeping. The \code{axi()}
constructor instantiates the five physical channels alongside the
intermediate protocol queues. It returns a composite \code{Link} containing
a \code{Host} endpoint, a \code{Per} endpoint, and the hardware ports for
both protocol units. Top-level designs wire the two protocol units together
like standard hardware modules while application processes manipulate the
endpoint handles.

\begin{figure}[htbp]
\centering
\begin{lstlisting}[style=ascii,frame=none,numbers=none]
  host client        issue / beat / grant / done / release
  AxiHost            a unit, lowered
     aw ar w b r     the AXI4 channels, where a router goes
  AxiPer             a unit, lowered
  peripheral client  request / write data / answer / read beat
\end{lstlisting}
\caption{Layers of an AXI link. The physical layer implements five standard
AXI4 channels, while client software and peripheral handlers operate on
discrete transactions.}
\label{fig:a-layers}
\end{figure}

\section{The Client}
\label{sec:a-client}

The host interface provides two primary methods. The \code{read} method
takes a target address and a beat count. The \code{write} method takes an
address and a slice of payload beats. Both return as soon as the initial
burst enters the channel, yielding a completion handle. Multiple handles
may remain outstanding simultaneously, and the host may await them in
arbitrary sequence.

\axipart{hostapi}{The host endpoint interface from
\code{lib/parts/src/bus/axi.rs}.}{lst:a-hostapi}

The length of a burst derives directly from the provided data slice,
preventing mismatches between declared length and submitted beats. The
\code{write} helper automatically asserts \code{last} on the final beat.
Hardware tracking logic allocates identifiers dynamically, keeping
identifier assignment transparent to the client. The client specifies only
addresses, data words, and strobes.

The peripheral interface exposes a single entry point. The \code{accept}
method blocks until a complete transaction arrives and yields a decoded
structure: either a read requesting a specific beat count, or a write with
all data beats collected. The transaction object provides completion
methods: \code{ok}, \code{err}, or \code{data}.

\axipart{accept}{The peripheral endpoint interface yielding decoded
transactions.}{lst:a-accept}

A read transaction becomes ready immediately upon receipt of its address
beat. A write transaction becomes ready only after its final data beat
arrives. Because AXI4 omits transaction identifiers on the write data
channel, beats correlate strictly with the oldest active write address.
Restricting active gathering to a single transaction at a time prevents
concurrent worker processes from corrupting each other's data streams.

\section{Accepting Many}
\label{sec:a-many}

High-throughput peripherals overlap transaction execution. A controller
accepts requests continuously, dispatches them to processing pipelines, and
posts responses as individual operations finish. Each accepted transaction
holds its unique identifier alongside a completion handle to the endpoint
response queue. Handlers may transfer this token across process boundaries
and complete operations out of order.

The \code{Open} table indexes pending transactions by identifier. The
\code{serve} helper orchestrates multi-process dispatch through a single
abstraction: it spawns a configured pool of workers, where each worker
accepts incoming work whenever idle and retires completions independently.

\axipart{serve}{\code{serve} helper orchestrating bus transactions across
concurrent workers.}{lst:a-serve}

The \code{ex\_axi\_serve} example implements equivalent peripherals through
manual process coordination and through \code{serve}, asserting that both
configurations produce identical results. The manual version deploys three
distinct processes: an ingestion loop forwarding work tickets, a pool of
execution workers, and a retirement process drawing completed entries from
\code{Open}. The retirement process interacts solely with completion handles
without touching the ingest queue.

\begin{figure*}[t]
\centering
\input{axi_serve_timing.tex}
\caption{Ten one-word reads over four identifiers, three in flight. An
identifier returns to the allocation pool on \code{release} only after the
client consumes the corresponding burst response. The \code{busy} vector
tracks active allocations.}
\label{fig:a-serve}
\end{figure*}

\lstinputlisting[caption={\code{ex\_axi\_serve.rs}. Run with
\code{bazel run //lib/examples:ex\_axi\_serve}.},label={lst:a-serve-ex}]{lib/examples/ex_axi_serve.rs}

\lstinputlisting[style=ascii,caption={What \code{ex\_axi\_serve}
printed when the build ran it: the two peripherals, each answering
in the order its work finished.},label={lst:a-serve-out}]{out_axi_serve.txt}

\section{The Beats}
\label{sec:a-beats}

Signal definitions match the full AXI4 specification. An address phase
provides \code{id}, \code{addr}, \code{len}, \code{size}, \code{burst},
\code{lock}, \code{cache}, \code{prot}, \code{qos}, and \code{region}.
Because \code{aw} and \code{ar} channels share identical field structures,
TxHDL unifies them under a common address struct.

\axipart{beats}{Channel beat definitions for AXI4.}{lst:a-beats}

Bit widths appear as explicit generic parameters across the library.
Parameter \code{A} defines address width, \code{D} defines data width, and
\code{S} defines the byte-strobe width (\(D/8\)). Parameter \code{I}
specifies identifier width, yielding \(2^I\) available identifiers via
\code{NIDS}. Computing strobe width directly as \code{U<\{D / 8\}>} requires
nightly Rust compiler features. The library relies on explicit parameters
to maintain compatibility with stable compilers, following the conventions
established by \code{Fifo} and \code{Station}.

\section{What This Link Orders}
\label{sec:a-order}

The AXI specification establishes no ordering guarantees between read and
write channels. Because read and write paths operate independently, a
subordinate device may retire transactions in either sequence. A master
enforces ordering between writes and subsequent accesses by awaiting the
write response before issuing the dependent request.

Ordering constraints become critical when processors execute posted stores.
Software preparing direct-memory buffers for peripheral masters, such as
video scanout or Ethernet controllers, must confirm data residency before
ringing peripheral doorbells. The frame payload reaches memory through one
router port while the activation signal targets a control register on
another port. The bus fabric provides no implicit synchronization between
them. The processor core addresses this via \code{fence} instructions:
execution tracking registers increment on posted stores and decrement on
received responses. A \code{fence} stalls the execution pipeline until the
unacknowledged store counter reaches zero.

Earlier device drivers enforce ordering by performing a loopback read of the
final buffer word prior to asserting peripheral doorbells. The read access
stalls the core until the memory subsystem returns data. This pattern
succeeds only if memory cannot return stale contents while a preceding store
to the same location remains pending.

Target hardware guarantees this ordering property by construction, verified
through empirical cycle traces:

\begin{center}\footnotesize
\begin{tabular}{@{}ll@{}}
\toprule
Store issued & tick 22 \\
Load completed & tick 38 \\
Store completed & tick 40 \\
\bottomrule
\end{tabular}
\end{center}

In this measurement, the load completes two clock ticks before the write
response arrives, demonstrating that the store remained active while the load
returned newly written data. The test
\code{a\_load\_behind\_an\_unanswered\_store\_to\_one\_address} in
\code{//lib/parts:parts\_test} enforces this behavior. The test asserts both
the cycle ordering and the resolved value, flagging an error if the store
finishes ahead of the load.

This dynamic yields two distinct conclusions.

First, drivers targeting this architecture can reliably use readbacks to
guarantee memory updates. This mechanism allows drivers to satisfy the
synchronization contracts documented in \code{EthSlots}.

Second, this guarantee reflects the specific implementation of this
interconnect fabric rather than universal AXI compliance. A fabric
implementation that retired read bursts ahead of pending write buffers
would invalidate driver assumptions, which the regression test suite
detects automatically.

\section{The Trackers}
\label{sec:a-trackers}

Two protocol engines manage channel interactions. Both modules reside in the
lowered hardware subset, compiling directly to netlists while participating
in software simulation.

\code{AxiHost} translates transaction requests into address beats on the
appropriate physical channel, allocates identifiers, forwards write
payloads, and routes inbound responses. Identifier allocation follows a
round-robin schedule across bursts.

\axipart{hostunit}{\code{AxiHost}: host-side hardware protocol engine.}{lst:a-hostunit}

An identifier remains reserved from its address phase until the client
explicitly releases it upon consuming the response. Releasing the tag at
consumption rather than on initial bus response prevents aliasing: recycling
an identifier immediately upon bus completion could reassign that tag while
prior data remains uncollected in the receive queue, causing response
collision. The allocation arbiter skips reserved slots and selects the next
available index, stalling transaction issue only when all \(2^I\) tags
remain active simultaneously. This mechanism provides natural backpressure.
An earlier design waited strictly for the next sequential identifier,
deadlocking whenever a client retained a single unconsumed handle; the link
fuzzer in Section~\ref{sec:a-fuzz} isolated that defect.

\code{AxiPer} arbitrates between read and write address channels to generate
a unified request stream for the peripheral. It forwards write payloads to
the target core and encodes peripheral acknowledgments into AXI response
beats. Address channels arbitrate via round-robin selection, preventing
continuous read traffic from starving write operations.

\axipart{perunit}{\code{AxiPer}: peripheral-side hardware protocol engine.}{lst:a-perunit}

\section{Supplying Your Own}
\label{sec:a-own}

The standard protocol engines maintain no privileged status. The
\code{axi()} constructor yields their interfaces as concrete port types:
\code{HostIn}, \code{HostOut}, \code{PerIn}, and \code{PerOut}. Custom units
implementing these typed ports can replace the default trackers directly.
Dedicated hardware peripherals can also bypass \code{AxiPer} entirely and
bind to the raw five-channel AXI bus, as demonstrated by the SDRAM
controller.

\section{A Client That Is Hardware}
\label{sec:a-hardware}

The \code{Host} and \code{Per} abstractions serve software simulation and
testbenches. \code{Host} buffers arriving completions to satisfy
asynchronous futures, while \code{Per} accumulates burst payloads into heap
vectors. The \code{Open} lookup table tracks in-flight transactions, and
\code{serve} spawns simulation threads on demand. These software constructs
do not lower to synthesizable hardware.

Synthesizable clients interact with the underlying hardware channels
directly. The host hardware interface bundles outbound channels (\code{issue},
write payload, and tag release) alongside inbound channels (allocation
grant, write response, and read data). The peripheral interface receives
request descriptors and write payloads while transmitting completion
responses and read data beats. The \code{axi\_units()} constructor
establishes this topology, returning raw \code{HostClient} and
\code{PerClient} bundles that synthesize cleanly with surrounding logic.

The companion constructor \code{axi\_to\_unit()} connects a software
testbench host to a synthesizable hardware peripheral:

\begin{lstlisting}[style=ascii,frame=none,numbers=none]
  axi()         -> Host / Per          simulation testbench pair
  axi_units()   -> HostClient /        synthesizable hardware pair
                   PerClient
  axi_to_unit() -> Host / PerClient    simulation host driving
                                       synthesizable peripheral
\end{lstlisting}

The rasterizer core in Razboj~\cite{gpu} uses \code{PerClient} as a
synthesizable bus endpoint. The underlying fabric remains identical across
configurations: the same protocol units, standard five-channel bus signals,
and identifier management rules apply throughout.

\section{The Simplest Peripheral}
\label{sec:a-simplest}

A single one-word register serves as the minimal hardware peripheral.
The example \code{ex\_axi\_reg} implements this module as a synthesizable
client. It binds the four channels of \code{PerClient}: reading request
descriptors and write beats, while driving completion responses and read
data beats. Listing~\ref{lst:a-reg} provides the hardware definition, and
Listing~\ref{lst:a-reg-ex} provides the test driver.

\lstinputlisting[rangebeginprefix=//\ begin\{,rangebeginsuffix=\},%
rangeendprefix=//\ end\{,rangeendsuffix=\},includerangemarker=false,%
linerange=register-register,caption={\code{Register}, the simplest
peripheral.},label={lst:a-reg}]{lib/examples/ex_axi_reg.rs}

During each clock cycle, the controller inspects the request queue. A read
request proceeds as soon as the outbound read-data channel indicates
available capacity. The unit returns the stored data word, the matching
transaction identifier, and an asserted \code{last} flag in the same cycle.
A write request proceeds when its corresponding data beat arrives and the
response queue can accept an acknowledgment. The internal register latches
the payload, and an acknowledgment posts immediately. Consuming write
descriptors and data beats concurrently eliminates intermediate buffering
registers, leaving the single 32-bit storage cell as the only sequential
state.

This peripheral focuses strictly on single-beat word transactions. The unit
omits address decoding, relying on upstream routers to direct matching
traffic. It assumes single-beat bursts; bursts exceeding one beat are
invalid. It ignores byte strobes, updating the entire word on every write.
The interrupt controller in the Vreteno processor~\cite{vreteno} extends
this pattern to five addressed registers with strobe support, while the
framebuffer in Razboj~\cite{gpu} expands it to full memory arrays with burst
sequencing.

The testbench drives the register using an \code{axi\_to\_unit()} link. The
host software writes two consecutive words, reads both locations back, and
logs simulation timestamps at two steps per clock cycle.

\begin{figure}[htbp]
\centering
\input{axi_reg_timing.tex}
\caption{The first write and its read. The write's request and its
beat arrive together on \code{req} and \code{wd}. A cycle later both are
accepted, \code{word} updates, and the response emits on \code{ans}.
The read completes on \code{rb} one cycle after its request arrives.}
\label{fig:a-reg}
\end{figure}

\lstinputlisting[rangebeginprefix=//\ begin\{,rangebeginsuffix=\},%
rangeendprefix=//\ end\{,rangeendsuffix=\},includerangemarker=false,%
linerange=main-main,caption={The run, the rest of
\code{lib/examples/ex\_axi\_reg.rs}. Run it with
\code{bazel run //lib/examples:ex\_axi\_reg}.},%
label={lst:a-reg-ex}]{lib/examples/ex_axi_reg.rs}

\lstinputlisting[style=ascii,caption={What \code{ex\_axi\_reg}
printed when the build ran it: each answer and its time step, and
the register's netlist.},label={lst:a-reg-out}]{out_axi_reg.txt}

The synthesized netlist contains one 32-bit register accompanied by
combinational decode logic. Continuous integration tests compile this
netlist into VHDL for \code{nvc} and Verilog for Verilator, verifying
cycle-accurate port activity against the simulation model.

\section{Peripherals That Only Simulate}
\label{sec:a-sim}

Development environments frequently require simulation-only peripherals.
Testbenches load executable firmware images, architectural models deposit
datasets, and early bring-up procedures require memory long before physical
controllers are synthesized. The module \code{bus::axi::sim} isolates these
simulation components, which remain distinct from synthesizable targets.

\simpart{ram}{A RAM behind a peripheral end, from
\code{lib/parts/src/bus/axi/sim.rs}.}{lst:a-ram}

Listing~\ref{lst:a-ram} defines a memory model backed by a standard heap
vector. Test code populates program binaries into this model before
execution and inspects resulting state after simulation completes.
Implementing the model in standard Rust avoids hardware synthesis overheads:
a megabyte occupies a native \code{Vec<u8>} rather than an expensive array
of synthesized registers. Burst handling, byte strobes, and boundary bounds
resolve in concise software logic. In-bounds transfers complete with
\code{Okay}, while out-of-bounds accesses return \code{SlvErr}.

Synthesizable peripherals follow the hardware path, compiling from the
lowered subset directly to netlists as shown in Razboj's
framebuffer~\cite{gpu}. These approaches complement each other across design
phases: verification starts with flexible simulation models before
transitioning to synthesized implementations, without modifying upstream
client interfaces.

\section{Peer to Peer}
\label{sec:a-peer}

The example \code{ex\_axi} demonstrates point-to-point communication between
one host and one peripheral over a single link. The host dispatches three
concurrent bursts (two reads and one write) without blocking on completions.
Each operation receives an independent transaction identifier. The
peripheral processes requests across three worker slots whose execution
latencies vary by target address, causing operations to complete out of order.

\begin{figure*}[t]
\centering
\input{axi_timing.tex}
\caption{The peer-to-peer run: three bursts issued, identifiers granted,
address beats driven on \code{ar} and \code{aw}, write beats marked with
\code{w\_last}, and completions returned out of order.}
\label{fig:a-peer}
\end{figure*}

\lstinputlisting[caption={\code{ex\_axi.rs}. Run with
\code{bazel run //lib/examples:ex\_axi}.},label={lst:a-peer-ex}]{lib/examples/ex_axi.rs}

\lstinputlisting[style=ascii,caption={What \code{ex\_axi} printed
when the build ran it: the three bursts, the order they were
answered in, and then the lowering, two modules.},label={lst:a-peer-out}]{out_axi.txt}

The host awaits transaction handles in their original submission order while
receiving the correct payload for each burst. Responses arriving out of
sequence accumulate in an internal endpoint inbox until the corresponding
future is polled. Polling any active handle drains the inbound response
channel, ensuring continuous queue throughput without requiring dedicated
background collector processes.

\section{The Netlist}
\label{sec:a-netlist}

Listing~\ref{lst:a-peer-out} concludes with the generated Verilog netlists
for both protocol trackers. Each channel interface expands into canonical
\code{\_data}, \code{\_valid}, and \code{\_ready} signals. For example,
\code{AxiPer} exposes \code{aw\_data}, \code{aw\_valid}, \code{aw\_ready},
and \code{ar\_data}. Field definitions pack into single bus vectors rather
than disjoint wires such as \code{awaddr} and \code{awlen}; unpacking them
requires only slice index definitions.

Sequential state comprises a four-bit allocation register, an active mask,
and a one-bit arbitration register. Step logic assigns clear names to
internal control nets: \code{free}, \code{go}, \code{wr\_go}, \code{rel\_go},
\code{take\_ar}, and \code{take\_aw}. The test target \code{ex\_axi} emits
both netlists during execution, and Bazel validates them against cycle
traces using \code{nvc} for VHDL and Verilator for Verilog.

Reserved language keywords require automatic sanitization. The internal
register storing the next scheduled identifier is named \code{next}, which
is a reserved keyword in VHDL. Rather than rejecting the identifier, the
compiler escapes it to \code{next\_rw} across both netlists and simulation
trace binders, resolving issue 497.

\section{Fuzzing the Link}
\label{sec:a-fuzz}

Targeted unit tests validate individual protocol invariants along
predetermined paths. In contrast, the randomized test suite
\code{bus::axi::fuzz} validates all invariants concurrently under randomized
stimuli. A pseudorandom host client dispatches bursts between 1 and 16 beats
across mixed read and write operations at arbitrary addresses, sustaining up
to 4 concurrent transactions and awaiting responses out of order. A companion
peripheral processes up to 4 concurrent transactions with randomized delays,
returning \code{SlvErr} for accesses within an unmapped test window.
Stalling relays instrument all 15 active channels (the 5 physical AXI buses
and the 10 internal client channels), injecting single-cycle stalls during
25\% of clock ticks and burst stalls lasting 8 to 47 cycles intermittently.

\fuzzpart{relay}{A relay, from
\code{lib/parts/src/bus/axi/fuzz.rs}.}{lst:a-relay}

The test harness enforces foundational protocol invariants. Each transaction
completes exactly once. Read operations return the exact data deposited by
preceding writes, guaranteed by avoiding address collisions with transactions
currently in flight. The host never reallocates an identifier while its
response remains uncollected. All test iterations complete within strict
cycle bounds, confirming absence of deadlocks. An autonomous wire monitor
inspects physical channel nets to verify that \code{last} asserts solely on
the terminal beat of every burst.

\fuzzpart{monitor}{The monitor of the AXI channels.}{lst:a-monitor}

To ensure comprehensive coverage across seeds, the fuzzer verifies that
corner cases actually trigger during execution. Test assertions verify that
completions and poll operations occur out of order, that identifiers wrap
around during pending operations, that maximum-length 16-beat bursts
execute, and that relays stall every channel for 8 cycles or more. The
standard regression suite in \code{//lib/parts:parts\_test} evaluates 24
seeds; \code{//lib/parts:axi\_fuzz\_test} runs dedicated sweeps; and
continuous integration evaluates 1024 seeds nightly using
\code{AXI\_FUZZ\_SEEDS}. A complementary unit test truncates a 4-beat read
response to 3 beats, verifying that the peripheral endpoint traps the defect
immediately rather than hanging the client indefinitely.

Early fuzzer runs exposed two subtle bugs. First, when concurrent worker
processes targeted the peripheral response port in the same cycle, one beat
overwrote the other (issue 182); arbitration logic now serializes concurrent
writes across consecutive cycles. Second, the host tracker originally stalled
on unavailable sequential identifiers even when alternative tags were free
(issue 184); the scheduler now rotates to the next available identifier
dynamically. Mutation tests that misplace the \code{last} flag trigger
assertions on the initial corrupt read.

\section{The Router}
\label{sec:a-routing}

The five physical channels connecting \code{AxiHost} and \code{AxiPer}
transport standard transaction packets. A router therefore functions as a
standard hardware unit rather than an ad-hoc protocol extension. It accepts
outbound \code{aw}, \code{ar}, and \code{w} channels from a host and
replicates them across an array of peripheral ports. In the reverse
direction, it concentrates per-peripheral \code{b} and \code{r} channels
into single upstream host interfaces.

Peripheral address windows are specified as base-and-mask pairs in a mapping
struct. Encoding the address map into the Rust type system produces fully
static decoders in the synthesized netlist:

\begin{lstlisting}
struct ThreeMap;

impl AddrMap<3> for ThreeMap {
    const RANGES: [(usize, usize); 3] =
        [(0x1000, 0xf000), (0x2000, 0xf000), (0x3000, 0xf000)];
}

type Rtr = Router<3, ThreeMap, 16, 32, 4, 3>;
\end{lstlisting}

A single generic module scales to arbitrary peripheral counts. Peripheral
connections are declared as port arrays, which the lowering stage unrolls
into dedicated wiring nets during synthesis. Address ranges evaluate
identically at elaboration time, following the \code{AddrMap} patterns
described in~\cite{examples}.

\subsection*{State Requirements}

A bus router enforces three critical protocol constraints beyond
straightforward signal demultiplexing.

First, write data beats omit transaction identifiers. The AXI4 standard
places identifiers on \code{aw} but omits them from \code{w}. Consequently,
write data beats associate implicitly with the oldest uncompleted address
phase. The router tracks this destination through state registers:
\code{wsel} stores one selection bit per peripheral, and \code{wbusy}
indicates an ongoing burst. Subsequent write address phases stall until the
final beat of the active burst transfers. Without this serialization,
interleaved data beats would corrupt destination buffers across distinct
peripherals.

Second, read data beats belonging to a burst must transfer contiguously.
Because host logic relies on \code{rlast} to detect burst termination,
peripheral arbitration locks on the initial read beat. The router records the
selected target in \code{rsel} and holds it via \code{rbusy} until the burst
terminates.

Third, unmapped target addresses require explicit protocol termination.
Dropping invalid requests would leave host handles waiting indefinitely,
starving the tag allocator. The router handles unmapped accesses internally:
returning \code{DecErr} on the \code{b} channel after consuming write data
beats, and generating \code{DecErr} beats across the \code{r} channel matching
the exact burst length requested in \code{rerr\_left}. Emitting the precise
requested beat count is required to complete the client receive loop
cleanly.

\subsection*{The Unit}

Listing~\ref{lst:a-rstate} presents the internal register state of the
router. Listing~\ref{lst:a-rdecode} details address decoding and peripheral
ready logic, while Listing~\ref{lst:a-rmerge} presents the upstream
multiplexer merging peripheral responses with local decode errors. The
lowering step unrolls port iterations into explicit named nets, documenting
precise inter-peripheral dependencies in the netlist.

\subsection*{Validation Trace}

The example \code{ex\_axi\_router} configures one host, three mapped
peripherals, and an unmapped address range. The host issues five bursts
concurrently: reads to the first two peripherals, a write to the third, and
both a read and a write targeting unmapped space. The peripherals process
requests with independent latencies, returning out-of-order completions
tagged with their respective peripheral indices as verified by test
assertions.

\begin{figure*}[t]
\centering
\input{axi_router_timing.tex}
\caption{The router's run: address phases arriving on \code{aw} and
\code{ar}, peripheral selections asserted, decode error flags raised by
unmapped accesses, and completions returned upstream with identifiers.}
\label{fig:a-router}
\end{figure*}

\lstinputlisting[caption={\code{ex\_axi\_router.rs}. Run with
\code{bazel run //lib/examples:ex\_axi\_router}.},%
label={lst:a-router-ex}]{lib/examples/ex_axi_router.rs}

\lstinputlisting[style=ascii,caption={What \code{ex\_axi\_router}
printed when the build ran it: five bursts, each answered by the
right range, and then the router's netlist.},%
label={lst:a-router-out}]{out_axi_router.txt}

The host and peripheral implementations remain completely decoupled from
routing logic. The host operates as the peer client from
Section~\ref{sec:a-peer}, while peripherals instantiate standard \code{serve}
worker pools. Neither module contains address ranges or routing parameters.
Inserting a router requires only channel reconnection: the router binds to
the peripheral-facing ports of the host link and the host-facing ports of
each peripheral link.

The router netlist lowers to Verilog and VHDL, passing regression verification
under \code{nvc} and Verilator against cycle-accurate traces. Upstream
arbitration among multiple competing hosts completes the bus interconnect
architecture.

\section{The Arbiter}
\label{sec:a-arbiter}

The arbiter mirrors router functionality by multiplexing multiple host
channels into a single downstream peripheral interface. While the router
decodes destination addresses to fan out requests, the arbiter resolves host
contention to merge streams. Pairing arbiters with routers creates full
multi-master crossbars, enabling direct memory access for secondary
controllers alongside processor cores.

The module \code{Arbiter<N, ..>} supports up to 8 host interfaces through port
arrays (\code{Arbiter2} through \code{Arbiter8}). It implements round-robin
scheduling by default, advancing the grant token past the winning port to
prevent host starvation. Setting the priority parameter activates
fixed-priority arbitration, giving precedence to lower-indexed hosts.

\subsection{Response Routing via Tag Extension}

Because downstream peripherals do not distinguish individual upstream
masters, response routing relies on transaction identifiers. AXI4 preserves
address-phase identifiers in completion responses without modification. The
arbiter exploits this property to record request provenance.

Upon granting host access, the arbiter prefixes the host port index onto the
upper bits of the outbound identifier. On the return path, it strips these
routing bits to restore original host tags while steering responses to the
corresponding master. The peripheral interface width \(J\) must satisfy
\(J \ge I + \lceil\log_2 N\rceil\), where \(I\) represents the master
identifier width and \(N\) the host count. Host clients retain private
identifier spaces while the peripheral operates over a widened global tag
space.

\subsection{Write Channel Serialization}

Because the write data channel provides no identifiers, data beats correlate
strictly with the oldest uncompleted write address phase. The interconnect
fabric must maintain write beat order explicitly.

The arbiter locks the write data path to the host whose write address phase
won arbitration until the terminal beat passes. Competing write address
requests stall until the active burst completes. In contrast, read paths
require no persistent locking because each beat holds an explicit identifier
and can interleave safely. Read arbitration operates cycle by cycle, whereas
write arbitration asserts an atomic burst lock.

In the verification run \code{ex\_axi\_arbiter}, two masters submit four-beat
write bursts simultaneously. The arbiter serializes the requests into two
non-interleaved bursts at the downstream memory interface, verifying correct
burst isolation.

\begin{figure*}[t]
\centering
\input{axi_arbiter_timing.tex}
\caption{The arbiter's run: both hosts submitting write address phases and
payloads, \code{wsel} indicating active host selection during \code{wbusy},
consecutive four-beat bursts transferring without interleaving, and write
responses routed by host prefix tags.}
\label{fig:a-arbiter}
\end{figure*}

\lstinputlisting[caption={\code{ex\_axi\_arbiter.rs}. Run with
\code{bazel run //lib/examples:ex\_axi\_arbiter}.},%
label={lst:a-arbiter-ex}]{lib/examples/ex_axi_arbiter.rs}

\lstinputlisting[style=ascii,caption={What \code{ex\_axi\_arbiter}
printed when the build ran it: each host's own answers, the two write
bursts as the memory was given them, and then the arbiter's
netlist.},%
label={lst:a-arbiter-out}]{out_axi_arbiter.txt}

Neither host client nor peripheral logic references arbiters, port indices,
or composite identifiers. Designers wire the arbiter directly between
host-facing links and peripheral interfaces without modifying client
endpoints.

The arbiter netlist compiles to Verilog and VHDL, matching simulation traces
under \code{nvc} and Verilator. Unit tests in \code{lib/parts/src/bus/arbiter.rs}
evaluate configurations with 2 and 4 hosts, asserting proper response
routing, non-interleaved write bursts, and balanced throughput under
saturation.

\section{AXI-Lite and the Bridge}
\label{sec:a-lite}

Simple peripherals rarely require the full complexity of AXI4. Serial ports
and configuration register banks transfer one word at a time, needing neither
burst lengths nor transaction identifiers. The AXI-Lite specification
simplifies the five standard channels into single-beat transactions without
out-of-order execution. Listing~\ref{lst:a-lite-beats} defines the AXI-Lite
beat structures from \code{bus::axi\_lite}. AXI-Lite peripherals require no
complex tracking state: read channels accept an address and return a data
word, while write channels consume an address and data word to return a
completion status.

\lstinputlisting[rangebeginprefix=//\ begin\{,rangebeginsuffix=\},%
rangeendprefix=//\ end\{,rangeendsuffix=\},includerangemarker=false,%
linerange=beats-beats,caption={The AXI-Lite beats. The response is
AXI4's \code{Resp}.},label={lst:a-lite-beats}]{lib/parts/src/bus/axi_lite.rs}

The constructor \code{axi\_lite()} creates these simplified channels and
yields endpoint handles matching the conventions of \code{axi\_units()}.

\subsection*{The Bridge Architecture}

Bridging AXI4 host fabrics to AXI-Lite peripherals requires protocol
translation. The \code{LiteBridge} unit multiplexes an inbound AXI4 interface
across \(N\) downstream AXI-Lite endpoints via port arrays that unroll during
netlist synthesis (issue 500). Static address maps specify peripheral
boundaries through compile-time constants (issue 593):

\begin{lstlisting}
pub struct TwoMap;

impl AddrMap<2> for TwoMap {
    const RANGES: [(usize, usize); 2] =
        [(0x1000, 0xf000), (0x2000, 0xf000)];
}

type Bridge = LiteBridge<2, TwoMap, 16, 32, 4, 2>;
\end{lstlisting}

On its upstream interface, the bridge accepts \code{aw}, \code{ar}, and
\code{w} channels and drives \code{b} and \code{r} channels. It occupies the
topological position of an \code{AxiPer} unit, attaching either directly to a
host tracker or behind a downstream router port.

The bridge enforces four invariant rules:

First, it serializes incoming bursts. Because AXI-Lite lacks identifiers,
the bridge blocks subsequent transactions until the active burst retires
completely. Simultaneous read and write requests arbitrate fairly via
round-robin selection.

Second, the bridge translates multi-beat bursts into sequential single-beat
AXI-Lite requests. The first beat targets the base address; subsequent beats
advance by \(1 \ll \text{size}\) bytes, whereas fixed bursts maintain a
constant address. Read responses forward upstream with the original
transaction identifier and an asserted \code{last} flag on the terminal
beat. Write responses accumulate into a single status beat, capturing the
initial error status if any beat faults. The bridge handles wrapping bursts
as standard incrementing bursts.

Third, address decoding maps each burst to a specific peripheral line in
\code{sel}, resolving overlapping matches by lowest numerical port index.

Fourth, requests targeting unmapped addresses complete locally within the
bridge. Write beats are consumed and discarded with \code{DecErr} responses,
while read requests return \code{DecErr} across the exact number of beats
specified by the host request.

\litepart{state}{The bridge's state, for any count of
peripherals.}{lst:a-lite-state}

\litepart{accept}{Taking a burst: the turn between the address
channels, the decode, and the address step.}{lst:a-lite-accept}

\litepart{requests}{A beat's request, to the peripheral the burst
decoded to.}{lst:a-lite-beats-out}

\litepart{answers}{The answer to a beat, from the peripheral or from
the bridge itself, going up.}{lst:a-lite-answers}

\litepart{drives}{The registers: a burst taken, a request sent, a beat
answered.}{lst:a-lite-drives}

\subsection*{Validation Run}

The verification harness \code{ex\_axi\_lite} places a two-port bridge behind
an AXI4 host, addressing two register banks (\code{Regs}) alongside an
unmapped hole. Each bank comprises a compact synthesizable AXI-Lite
peripheral. The host executes a three-word write burst, a four-word read
burst, a three-beat fixed read burst, access cycles to the secondary bank,
and invalid access attempts to the unmapped window. Assertions confirm all
payload words and status flags. The bridge and peripheral banks lower to
Verilog and VHDL, passing co-simulation under \code{nvc} and Verilator.

\begin{figure*}[t]
\centering
\input{axi_lite_timing.tex}
\caption{The bridge's run: bursts arriving on \code{aw} and \code{ar}, bridge
busy state tracked in \code{sel}, remaining beat counts decrementing, and
individual AXI-Lite transactions dispatched to destination banks.}
\label{fig:a-lite}
\end{figure*}

\lstinputlisting[caption={\code{ex\_axi\_lite.rs}. Run with
\code{bazel run //lib/examples:ex\_axi\_lite}.},%
label={lst:a-lite-ex}]{lib/examples/ex_axi_lite.rs}

\lstinputlisting[style=ascii,caption={What \code{ex\_axi\_lite}
printed when the build ran it: each burst, the tick its answer
arrived, and the answer.},%
label={lst:a-lite-out}]{out_axi_lite.txt}

The Vreteno processor~\cite{vreteno} uses this bridge to drive its
memory-mapped serial interface. Placing the UART behind
\code{LiteBridge<1, ..>} eliminated three dedicated staging registers
previously needed to synchronize write data beats.

\section{Reaching Wishbone}
\label{sec:a-wb}

External intellectual property blocks often interface via Wishbone rather
than AXI. The DDR3 memory controller used on target hardware provides a
pipelined Wishbone interface. The bridge \code{bus::wb::AxiWb} bridges these
domains: it instantiates as a synthesizable peripheral client, terminating
AXI channel endpoints on one side and driving master lines for a pipelined
Wishbone bus on the other.

\lstinputlisting[rangebeginprefix=//\ begin\{,rangebeginsuffix=\},%
rangeendprefix=//\ end\{,rangeendsuffix=\},includerangemarker=false,%
linerange=state-state,caption={The bridge's state, from
\code{lib/parts/src/bus/wb.rs}.},%
label={lst:a-wbstate}]{lib/parts/src/bus/wb.rs}

The bridge processes bursts sequentially, converting multi-beat transactions
into individual Wishbone accesses. Read operations dispatch to the Wishbone
bus immediately; write operations dispatch once data beats arrive from the
AXI channel. The bridge asserts \code{cyc} and \code{stb} until the target
releases \code{stall}, indicating acceptance. It then maintains \code{cyc}
until the slave asserts \code{ack}, returning read data or write completion
back to the AXI fabric when channel buffering permits. All outbound Wishbone
signals register directly in output flip-flops.

\lstinputlisting[rangebeginprefix=//\ begin\{,rangebeginsuffix=\},%
rangeendprefix=//\ end\{,rangeendsuffix=\},includerangemarker=false,%
linerange=run-run,caption={The bridge, a cycle at a
time.},label={lst:a-wbrun}]{lib/parts/src/bus/wb.rs}

Wishbone uses word-aligned addresses. The bridge converts byte addresses to
word addresses by shifting right by two bits and truncating to target bus
width. Peripherals mapped within aligned address windows receive relative
offsets directly without runtime subtraction logic.

For verification, \code{bus::wb::sim::WbMem} provides an architectural
Wishbone memory model in native Rust. The model simulates controller
initialization by injecting an initial calibration stall of 10 cycles,
followed by a fixed 2-cycle latency per request.

\simwbpart{mem}{The Wishbone memory, from
\code{lib/parts/src/bus/wb/sim.rs}.}{lst:a-wbmem}

The example \code{ex\_axi\_wb} writes three words through the bridge and reads
them back. An \code{axi\_to\_unit()} construct couples the software host
testbench to the synthesizable bridge unit. The simulation engine schedules
memory model execution ahead of the bridge, allowing the bridge to sample bus
lines synchronously within the active clock cycle.

\begin{figure*}[t]
\centering
\input{axi_wb_timing.tex}
\caption{The first write through the bridge: request and data beat arrival,
\code{stb} asserted across memory calibration until accepted without stall,
\code{cyc} held until acknowledgment, and write response delivered on
\code{ans}.}
\label{fig:a-wb}
\end{figure*}

\lstinputlisting[caption={\code{ex\_axi\_wb.rs}. Run with
\code{bazel run //lib/examples:ex\_axi\_wb}.},%
label={lst:a-wb-ex}]{lib/examples/ex_axi_wb.rs}

\lstinputlisting[style=ascii,caption={What \code{ex\_axi\_wb} printed
when the build ran it: the lines in each cycle a request is on them,
the words read back, and the bridge's netlist.},%
label={lst:a-wb-out}]{out_axi_wb.txt}

The bridge is lowered, and the build simulates its netlist against
this run under \code{nvc} and under Verilator at its ports.

\begin{thebibliography}{9}

\bibitem{tutorial}
F.~Filmar and D.~Janković, ``TxHDL, step by step,'' project
repository \code{HDL/txhdl}, \code{//docs:tutorial}, 2026.

\bibitem{runtime}
F.~Filmar and D.~Janković, ``The TxHDL runtime,'' project repository
\code{HDL/txhdl}, \code{//docs:runtime}, 2026.

\bibitem{examples}
F.~Filmar and D.~Janković, ``TxHDL by example,'' project repository
\code{HDL/txhdl}, \code{//docs:examples}, 2026.

\bibitem{station}
F.~Filmar and D.~Janković, ``A reservation station in TxHDL,''
project repository \code{HDL/txhdl}, \code{//docs:station}, 2026.

\bibitem{gpu}
F.~Filmar and D.~Janković, ``Razboj: A minimal GPU in TxHDL,'' project
repository \code{HDL/txhdl}, \code{//docs:razboj}, 2026.

\bibitem{vreteno}
F.~Filmar and D.~Janković, ``Vreteno: an RV32IMAC core in TxHDL,''
project repository \code{HDL/txhdl}, \code{//docs:vreteno}, 2026.

\bibitem{cover}
F.~Filmar and D.~Janković, ``TxHDL documents,'' project repository
\code{HDL/txhdl}, \code{//docs:cover}, 2026.

\end{thebibliography}

\end{document}
